\chapter{Dopamine thực sự nói gì}
% full version: chapters/07_dofaalt.tex

Ở chương trước dopamine rất đơn giản: một con số, ``tốt hơn hay tệ hơn mong đợi''. Đó là phép gần đúng bậc nhất, và nó hoạt động tốt. Nhưng các thí nghiệm những năm gần đây cho thấy trên ``dây dẫn'' này có nhiều hơn một con số. Chương này nói về những chỗ bức tranh phức tạp hơn và những chỗ các nhà khoa học còn đang tranh luận.

\section*{Người lạc quan và người bi quan}

Các nơ-ron dopamine không giống nhau. Có nơ-ron phản ứng với bất ngờ dễ chịu mạnh hơn với bất ngờ khó chịu --- đó là những người lạc quan. Có nơ-ron thì ngược lại --- những người bi quan. Nếu mỗi nơ-ron học theo quy tắc riêng, người lạc quan sẽ học được phần thưởng trong trường hợp tốt nhất, còn người bi quan --- trong trường hợp tồi nhất. Cùng nhau, chúng ghi lại không chỉ một kỳ vọng trung bình mà cả dải kết cục có thể xảy ra, như nhà cái biết không chỉ ứng viên sáng giá nhất mà cả cơ hội của từng đối thủ. Ý tưởng này khớp với thí nghiệm; còn bộ não cần cả dải ấy để làm gì thì hiện vẫn là giả thuyết.

\pencilfig{7}{\pencilart{7}}{Mỗi nơ-ron dopamine giữ một điểm riêng trên đường cong các phần thưởng có thể: cùng nhau, chúng nhớ cả đường cong.}

\section*{Không chỉ ``tốt hơn'' mà còn ``quan trọng''}

Một phần nơ-ron dopamine cũng lóe lên với những sự kiện khó chịu nhưng quan trọng: tiếng động lớn, nguy hiểm. Có vẻ như trên mạng dopamine không chỉ có dấu của sự bất ngờ mà còn có tầm quan trọng của nó --- tín hiệu ``hãy chú ý''.

\section*{Hai tín hiệu trên một dây dẫn}

Những đỉnh vọt nhanh thì dạy. Còn mức dopamine trung bình chậm, thay đổi trong vài phút, có vẻ nói về chuyện khác: môi trường lúc này giàu có đến đâu. Nếu xung quanh nhiều phần thưởng, thời gian là quý --- và con vật hành động nhanh hơn. Thí nghiệm xác nhận rằng mức dopamine liên quan đến độ hăng hái của hành động. Kỹ sư sẽ gọi đây là tách tín hiệu theo tần số: tần số cao mang một thứ, tần số thấp mang thứ khác.

\section*{Dâng lên gần đích}

Khi con vật tiến gần tới phần thưởng, dopamine thường không bùng lên mà tăng dần đều. Một cách giải thích: đó chính là sự kỳ vọng, tín hiệu ``đích gần rồi, cố lên''. Cách khác giữ nguyên ý tưởng bất ngờ: từ xa con vật đánh giá tình hình mơ hồ, càng tới gần càng chính xác, và mỗi lần chính xác hơn là một bất ngờ dễ chịu nhỏ. Một thí nghiệm tinh tế trong hành lang ảo, nơi con vật bị chuyển tức thời tới gần đích hơn, ủng hộ cách thứ hai: dopamine đáp lại bằng một đỉnh vọt, đúng như một sự bất ngờ phải làm.

\section*{Nhìn lại phía sau}

Còn có một phương án táo bạo: dopamine không trả lời câu hỏi ``điều gì sẽ xảy ra?'' mà trả lời ``điều gì đã gây ra phần thưởng?'' --- nó nhìn về quá khứ chứ không phải tương lai. Những thí nghiệm mà hai lý thuyết dự đoán khác nhau hiện vẫn mâu thuẫn với nhau. Đây là một cuộc tranh luận khoa học đang diễn ra.

\section*{Vectơ thay cho con số}

Cuối cùng, dopamine còn phản ứng khi \emph{vị} của phần thưởng thay đổi mà giá trị của nó thì không, và giúp học cả khi không có phần thưởng nào. Có vẻ đây không phải một sai số mà cả một bó sai số --- mỗi sai số cho một đặc điểm của điều đã xảy ra.

Kết luận của chương: phần lớn các lý thuyết mới không phủ nhận bức tranh đơn giản mà mở rộng nó. Dopamine không phải một dây dẫn, mà là một bó dây nhỏ gồm vài sợi. Chỉ có ``nhìn lại phía sau'' là thực sự tranh cãi với nó.

\fullversion{7}
